\subsubsection{Frontend and Backend Development}

As a full-stack web application, the implementation is divided into two main areas. The frontend is responsible for the user interface (UI) and user experience (UX). The backend implements the complex implementations hidden in the frontend and integrates external services via a RESTful API.

\textbf{FBD1. Expense Uploads}

Create an intuitive visual interface for receipt submission via drag-and-drop or file upload. The backend will later handle database updates to register the new ticket. In this way, uploaded expenses will be accessible on future logins.

\textbf{FBD2. Uploaded Expenses Real-time Dashboard}

Call the backend API from the frontend app to recover all uploaded receipts of the authenticated user. The frontend dashboard displays each expense and its real-time processing status, which can be:

\begin{itemize}[leftmargin=0.5cm,noitemsep]
    \item \textbf{Processing:} The expense is pending processing by the AI model.

    \item \textbf{Pending Review:} The AI model has successfully analyzed the expense, the uploader can now review it.

    \item \textbf{Completed:} The user has reviewed and completed all required information. The expense is ready for being included in a new or existing report.

    \item \textbf{Failed:} The AI couldn't extract data from the receipt (due to unavailability or unrecognized format).

\end{itemize}

\newpage 

\textbf{FBD3. AI-Extracted Expense Data Review and Correction}

Let employees correct the AI-extracted expense data and fill in the missing required information in the frontend UI. Once the user saves the new expense information, the backend will then update the database with the corrected data.

\textbf{FBD4. Uploaded Expenses Deletion}

Allow employees to delete any uploaded receipt, regardless of its current status. The frontend provides visual tools to delete one or multiple expenses. After the user confirms the deletion, the backend marks them as deleted in the database.

\textbf{FBD5. Expense Reports View, Creation, Edit, Submission, and Deletion}

Allow employees to create reports from uploaded, AI-processed, and reviewed expenses. They can edit the report details, submit for approval, or delete unsubmitted reports using the frontend UI. The backend updates the database status based on user actions.


\textbf{FBD6. Company Email User Authentication}

Grant employees login access using their company email account through the login page. The backend, which is integrated with Microsoft Entra ID,  verifies the user access. On first login, the backend creates a new user profile; on return visits, it retrieves the user's relevant data (e.g., uploaded expenses).

\textbf{FBD7.  Role-based User Access}

Assign the correct  role to the logged-in user, either regular employee or administrator. Administrators have additional features unavailable to regular employees. 

When a user logs in, the backend retrieves their corresponding role, enabling the frontend to display only the functionalities specific to that role.


\textbf{FBD8. Submitted Report Approval and Rejection}

Enable administrators to view reports submitted by their supervised employees for approval or rejection of reimbursement budgets. The frontend handles the interface, while the backend updates the database accordingly.

\textbf{FBD9. Employee Spending Limit Configuration}

Provide administrative users the ability to set spending limits for employees under their supervision through the frontend UI. The backend is accountable for registering the new employee expense limit.

\textbf{FBD10. Submitted  and Reviewed Expenses Records Visualization}

Enable users to view their submitted report record with their corresponding current status. Administrators can also see their reviewed report record, including their approval or rejection decision. To do that, the frontend calls the backend API to retrieve the accurate information and displays it to the user through a visual dashboard.

\textbf{FBD11. Email Notification Service Integration}

Send email notifications to approvers for new pending reports, and to employees upon report approval or rejection. The backend handles all email notifications triggered by frontend actions.
